\chapter{شبكات عصبية تستطيع أن تتعلم: نضيف \nt{DOPA}}
\chaptermark{شبكات تتعلم}
\label{sec:dofa}
\begin{chaptergoals}
\item لماذا لا يجوز للمشبك أن يستعمل إلا إشارات محلية، ولماذا يستبعد ذلك الانتشار العكسي؛
\item ماذا تتعلم قاعدة هيب: الاتجاه الذي تتغير فيه المداخل أكثر من غيره؛
\item كيف يحوّل أثر الأهلية مع إشارة البثّ $\delta$ الضجيجَ إلى خطوات صعودًا على تدرج المكافأة؛
\item كيف يعطي الناقد $\delta$ في الزمن المتصل، ولماذا يتعلم السلك المشترك الواحد ببطء.
\end{chaptergoals}

\section{قواعد اللعبة}

في كل دارات الفصول السابقة كان المهندس هو من يضع الأوزان. وفي الكائن الحي لا يوجد مهندس. يجب أن تستقر الأوزان وحدها، أثناء العمل، وبحيث تفعل الشبكة شيئًا مفيدًا. وهذا ما يسمى التعلّم.

إلى القواعد السابقة تضاف قاعدة أخرى. المشبك يغيّر وزنه بنفسه، ولا يستطيع أن يعرف إلا ما يجري \emph{بقربه}: نشاط المرسِل $x$، ونشاط المتلقي $y$، والمواد التي تصل إليه. لا أحد يستطيع أن يرسل إلى المشبك تصحيحًا خاصًا به.

وهذا يستبعد الطريقة الرئيسية في الشبكات العصبية الاصطناعية: الانتشار العكسي للخطأ. هناك يعبر خطأ الخرج الشبكة في الاتجاه المعاكس عبر الأوزان نفسها، في طور مستقل بعد التمرير الأمامي، ويتلقى كل وزن تصحيحه. ولهذا تلزم أسلاك عكسية تكرر الأمامية بدقة، ونبضة ساعة مشتركة. ولم يُعثر في الدماغ على أي منهما، على الأقل بالشكل الذي هما عليه في الشبكات الاصطناعية~\cite{Lillicrap2020,WhittingtonBogacz2019}.

سنكتب تغيّر الوزن عملية متصلة بطيئة:
\begin{equation}
\frac{\dd w}{\dd t} \;=\; \eta\,\Phi(\text{ما يعرفه المشبك}),
\label{eq:plasticity}
\end{equation}
حيث $\eta$ معدل التعلّم، وهو صغير إلى حد أن الوزن يكاد لا يتغير خلال نبضة واحدة. والفصل كله عن الشكل الذي يمكن أن تتخذه الدالة $\Phi$.

\section{أين يُخزَّن الوزن}
\label{sec:weightstore}

ما المشبك الذي ”يغيّر وزنه“؟ للوصلة جانبان: نهاية سلك المرسِل (الجانب \emph{قبل المشبكي}) وقطعة من غشاء المتلقي عليها المستقبلات (الجانب \emph{بعد المشبكي})، وبينهما شق ضيق. في وصلات \nt{GLUT} يقع غشاء المتلقي عادةً على \emph{شويكة}، وهي نتوء صغير على فرع الدخل عند المتلقي. من المناسب أن نفكك وزن الوصلة إلى ثلاثة عوامل~\cite{delCastillo1954}:
\begin{empheq}[box=\keybox]{equation}
w \;=\; N_s\,p\,A_1 .
\label{eq:quantal}
\end{empheq}
هنا $N_s$ عدد مواقع التحرير بين العصبونين، و$p$ احتمال أن تحرّر النبضة دفعة من الناقل في موقع واحد (وهي $p$ نفسها من الفصل~\ref{sec:code})، و$A_1$ استجابة المتلقي لدفعة واحدة، أي عدد المستقبلات التي تلتقطها. العامل $p$ يقيم عند المرسِل، و$A_1$ عند المتلقي، أما $N_s$ فيتطلب الجانبين: مواقع تحرير جديدة تنمو، وقديمة تختفي، ويستمر هذا طوال الحياة.

\emph{القرار للمتلقي.} في مشبك \nt{GLUT} النموذجي لا يمرّر أحد مستقبلات الغلوتامات التيار إلا إذا اجتمع شرطان: وصل الغلوتامات من المرسِل، وغشاء المتلقي مستثار أصلًا~\cite{Nowak1984}. هذه قاعدة هيب مدمجة في جزيئة واحدة، أي كاشف التزامن من الفصل~\ref{sec:glutcomputer} وقد وُضع على مدخل المشبك نفسه. وحين يعمل المستقبل يُدخل الكالسيوم إلى المتلقي، فيطلق الكالسيوم تغيّر الوزن. المتلقي هو الموضع الوحيد الذي يُرى فيه الدخل والخرج معًا، لذلك يُتخذ القرار عنده.

\emph{التنفيذ يمكن أن يتولاه أي جانب.} أكثر أنواع التقوية الطويلة الأمد دراسةً تجري عند المتلقي: تُدمج في الغشاء مستقبلات جديدة~\cite{Malinow2002}، وتكبر الشويكة~\cite{Matsuzaki2004}، أي يكبر $A_1$. لكن المتلقي يستطيع أن يغيّر $p$ أيضًا: فهو يحرّر مادة تعبر الشق عائدة إلى المرسِل، وهي القناة العكسية من الملحق~\ref{sec:othermediators}، فيبدأ المرسِل بتحرير ناقل أقل~\cite{WilsonNicoll2001}. أما التغيرات القصيرة الأمد في الوزن، أي الاستنزاف والتيسير من الفصل~\ref{sec:code}، فهي تغيرات في $p$ يديرها المرسِل وحده بحسب نشاطه القريب.

بالنسبة إلى المهندس، سجل الوزن مقسوم بين رقاقتين. قاعدة التعلّم ينفذها المتلقي، ويستطيع أن يكتب النتيجة عنده أو عند المرسِل عبر القناة العكسية. لذلك فكلمة ”محلي“ في قواعد هذا الفصل لا تعني جانبًا واحدًا من الوصلة، بل الوصلة كلها.

\section{التعلّم بلا معلّم}

أبسط ما يمكن أن يفعله المشبك أن يقوى حين يكون المرسِل والمتلقي نشطين معًا:
\begin{equation}
\frac{\dd w}{\dd t} \;=\; \eta\,x\,y .
\label{eq:hebb}
\end{equation}
هذه قاعدة هيب~\cite{Hebb1949}. إنها محلية ولا تحتاج إلى ساعة. لكن فيها العيب نفسه الذي في شبكة \nt{GLUT} في الفصل~\ref{sec:glutcomputer}: التغذية الراجعة الموجبة. الوزن القوي يكبّر $y$، و$y$ الكبير يقوّي الوزن أكثر، فتنمو الأوزان بلا حد. والعلاج تغذية راجعة سالبة على الوزن: ليتناقص الوزن أيضًا بما يتناسب مع $y^2$. لعصبون مداخله $\mathbf x$ وخرجه خطي $y=\mathbf w\cdot\mathbf x$ نحصل على
\begin{empheq}[box=\keybox]{equation}
\frac{\dd\mathbf w}{\dd t} \;=\; \eta\,y\,\bigl(\mathbf x - y\,\mathbf w\bigr) .
\label{eq:oja}
\end{empheq}
ما زالت القاعدة محلية: يعرف المشبك دخله $x_i$ والخرج $y$ ووزنه $w_i$.

\begin{keyblock}
\begin{proposition}[ماذا تجد قاعدة هيب]
\label{prop:oja}
تقود القاعدة~\eqref{eq:oja} متجه الأوزان إلى طول الوحدة على امتداد الاتجاه الذي تتغير فيه المداخل أكثر من غيره، أي إلى المتجه الذاتي الرئيسي لمصفوفة ارتباطات المداخل $C=\langle\mathbf x\mathbf x^{\mathsf T}\rangle$~\cite{Oja1982}. يتعلم العصبون أن يخرج أكثر مقدار واحد تعبيرًا يمكن أن تُضغط فيه مداخله.
\end{proposition}
\end{keyblock}

البرهان. نأخذ متوسط~\eqref{eq:oja} على المداخل: $\dd\mathbf w/\dd t=\eta\bigl(C\mathbf w-(\mathbf w^{\mathsf T}C\mathbf w)\,\mathbf w\bigr)$. النقاط الثابتة هي المتجهات الذاتية لـ$C$ ذات طول الوحدة. إذا كان $\mathbf w$ قريبًا من متجه ذاتي $\mathbf e_k$ قيمته الذاتية $\lambda_k$، فإن إضافة صغيرة على امتداد متجه آخر $\mathbf e_j$ تنمو مثل $e^{\eta(\lambda_j-\lambda_k)t}$. ولا تستقر إلا النقطة التي تكون عندها كل $\lambda_j<\lambda_k$، أي المتجه الرئيسي.

وثمة صيغة من قاعدة هيب تنظر إلى توقيت النبضات. إذا وصلت نبضة المرسِل \emph{قبل} نبضة المتلقي بـ$s$ ميلي ثانية، ازداد الوزن بمقدار $A_+e^{-s/\tau_+}$؛ وإذا وصلت \emph{بعدها} بـ$s$ ميلي ثانية، نقص بمقدار $A_-e^{-s/\tau_-}$، حيث $\tau_\pm$ من رتبة عشرين ميلي ثانية. هذه القاعدة مقيسة على مشابك عصبونات \nt{GLUT}~\cite{BiPoo1998}. إنها تقوّي الوصلات التي \emph{يمكن أن تكون سببًا} في الاستجابة، وتضعف المتأخرة (الشكل~\ref{fig:stdp}). مرّر عبر الشبكة سلسلة الإثارات نفسها مرات كثيرة، فتنمو فيها وحدها وصلات من كل حلقة إلى التي تليها: السلسلة من الفصل~\ref{sec:glutcomputer} وقد جُمعت بلا مهندس.

\begin{figure}[t]
\centering
\begin{tikzpicture}[>=Stealth,x=0.07cm,y=1.3cm]
\draw[->,gray!70!black] (-66,0) -- (68,0) node[right,black] {\small $s$};
\draw[->,gray!70!black] (0,-1.2) -- (0,1.35) node[above,black] {\small تغيّر الوزن};
\draw[very thick,blue!60!black] plot[domain=0.5:60,samples=50] (\x,{exp(-\x/20)});
\draw[very thick,red!70!black] plot[domain=-60:-0.5,samples=50] (\x,{-0.6*exp(\x/20)});
\draw[blue!60!black,thick] (0.5,0) -- (0.5,1);
\draw[red!70!black,thick] (-0.5,0) -- (-0.5,-0.6);
\foreach \x in {-40,-20,20,40}{\draw[gray!60!black] (\x,-0.04) -- (\x,0.04);}
\node[below,font=\scriptsize] at (20,-0.04) {$20\ms$};
\node[above,font=\scriptsize] at (-20,0.04) {$-20\ms$};
\node[align=left,font=\small,blue!60!black,anchor=west] at (22,0.95) {المرسِل أسبق:\\ الوصلة تقوى};
\node[align=right,font=\small,red!70!black,anchor=east] at (-12,-0.95) {المرسِل متأخر:\\ الوصلة تضعف};
\end{tikzpicture}
\caption{قاعدة هيب بحسب التوقيت. على المحور الأفقي $s=t_{\text{post}}-t_{\text{pre}}$، أي بكم تأخرت نبضة المتلقي عن نبضة المرسِل؛ $\tau_\pm=20\ms$، $A_-=0.6A_+$. عرض النافذة بضع عشرات من الميلي ثواني، أي من رتبة نافذة التزامن~\eqref{eq:window} نفسها.}
\label{fig:stdp}
\end{figure}

لكن كل قواعد هذا القسم تعلّم ما هو \emph{متكرر}، لا ما هو \emph{مفيد}. المشبك الذي لا يعرف إلا $x$ و$y$ لا يدري أجلب الفعلُ عصيرًا أم ألمًا. ولكي يتعلم المشبك من النتيجة يحتاج إلى إشارة ثالثة.

\section{العامل الثالث}

الإشارة الثالثة يأتي بها \nt{DOPA}: عصبونات الدوبامين تبثّ عددًا واحدًا، هو خطأ التنبؤ بالمكافأة $\delta$~\eqref{eq:rpe} (الفصل~\ref{sec:neurontypes}). ليكن تغيّر الوزن متناسبًا مع جداء ثلاثة عوامل: نشاط المرسِل، ونشاط المتلقي، و$\delta$.

الصعوبة أن المكافأة لا تأتي حين كانت الشبكة تختار الفعل، بل بعد مئات الميلي ثواني أو ثوانٍ. وحين تصل $\delta$ يكون الجداء $xy$ قد صار صفرًا منذ زمن، ويحتاج المشبك إلى ذاكرة تفيد أنه شارك. وأبسط ذاكرة هي دارة \foreignlanguage{english}{RC} التي نعرفها:
\begin{empheq}[box=\keybox]{equation}
\tau_e\,\frac{\dd e}{\dd t} \;=\; -e \;+\; x\,y ,
\qquad
\frac{\dd w}{\dd t} \;=\; \eta\,\delta(t)\,e(t) .
\label{eq:threefactor}
\end{empheq}
يسمى المقدار $e$ \emph{أثر الأهلية}~\cite{Fremaux2016}. يترك النشاط المشترك في المشبك علامة تذوب خلال الزمن $\tau_e$. إذا وصل الدوبامين خلال هذا الزمن تغيّر الوزن بإشارة $\delta$؛ وإلا اختفت العلامة بلا أثر (الشكل~\ref{fig:trace}). وقد قيس على مشابك \nt{GLUT} أن الدوبامين الذي يصل خلال ثانية أو ثانيتين بعد النشاط المشترك يحوّله إلى تقوية للوصلة، أما الذي يصل بعد ذلك فلا~\cite{Yagishita2014}. ووُجدت نوافذ لدونة طولها ثوانٍ أيضًا حيث لا تكون الإشارة الثالثة الدوبامين، بل الإثارة القوية للمتلقي نفسه~\cite{Bittner2017}.

\begin{figure}[t]
\centering
\begin{tikzpicture}[>=Stealth,x=7cm,y=1cm]
\fill[orange!12] (0.7,-0.2) rectangle (0.8,5.2);
% rows: x*y, delta, traces, weights
\foreach \y/\lab in {4.4/{$x\,y$},3.1/{$\delta$},1.4/{الأثر $e$},0/{الوزن $w$}}{
  \draw[gray!70!black] (0,\y) -- (1.42,\y);
  \node[left,font=\small] at (-0.02,\y+0.3) {\lab};}
\draw[very thick,blue!60!black] (0,4.4) -- (0,4.9) -- (0.2,4.9) -- (0.2,4.4);
\node[right,font=\scriptsize,blue!60!black] at (0.21,4.8) {المرسِل والمتلقي نشطان معًا};
\draw[very thick,orange!85!black] (0,3.1) -- (0.7,3.1) -- (0.7,3.7) -- (0.8,3.7) -- (0.8,3.1) -- (1.4,3.1);
\node[right,font=\scriptsize,orange!85!black] at (0.81,3.6) {وصل الدوبامين};
% traces, each in units of its own maximum
\draw[very thick,blue!60!black] plot[domain=0:0.2,samples=20] (\x,{1.4+1.1*(1-exp(-\x))/(1-exp(-0.2))})
  plot[domain=0.2:1.4,samples=60] (\x,{1.4+1.1*exp(-(\x-0.2))});
\draw[thick,densely dashed,gray!50!black] plot[domain=0:0.2,samples=30] (\x,{1.4+1.1*(1-exp(-\x/0.05))/(1-exp(-4))})
  plot[domain=0.2:1.4,samples=80] (\x,{1.4+1.1*exp(-(\x-0.2)/0.05)});
\node[right,font=\scriptsize,blue!60!black] at (1.0,2.15) {$\tau_e=1\,\text{s}$};
\node[right,font=\scriptsize,gray!40!black] at (0.3,1.65) {$\tau_e=50\ms$};
% weights
\draw[very thick,blue!60!black] (0,0.25) -- (0.7,0.25) -- (0.8,0.8) -- (1.4,0.8) node[right,font=\scriptsize] {ازداد الوزن عند $\tau_e=1\,\text{s}$};
\draw[thick,densely dashed,gray!50!black] (0,0.2) -- (1.4,0.2) node[right,font=\scriptsize,gray!40!black] {لا تغيير عند $\tau_e=50\ms$};
% time axis marks
\foreach \x/\l in {0/0,0.2/0.2,0.7/0.7,1.4/1.4}{\draw[gray!60!black] (\x,-0.05) -- (\x,0.05) node[below=3pt,font=\scriptsize] {$\l\,\text{s}$};}
\draw[<->,thick,gray!50!black] (0.2,5.35) -- node[above,font=\scriptsize,black] {تأخر المكافأة $0.5\,\text{s}$} (0.7,5.35);
\end{tikzpicture}
\caption{أثر الأهلية وفق القاعدة~\eqref{eq:threefactor}. يدوم النشاط المشترك $0.2\,\text{s}$، ويصل الدوبامين بعد نصف ثانية من انتهائه. الآثار معروضة كنسب من قيمتها العظمى. الأثر البطيء ($\tau_e=1\,\text{s}$) ما زال حيًا في تلك اللحظة، أما السريع ($\tau_e=50\ms$) فقد ذاب منذ زمن.}
\label{fig:trace}
\end{figure}

\begin{keyblock}
\begin{proposition}[التعلّم بإشارة بثّ]
\label{prop:threefactor}
لا تغيّر القاعدة~\eqref{eq:threefactor} إلا المشابك التي شاركت في عمل الشبكة خلال آخر $\tau_e$، وفي الاتجاه الذي تحدده الإشارة المشتركة $\delta$. إشارة البثّ مع الأثر المحلي تعطي تصحيحًا موجّهًا إلى عنوان. والتأخر بين الفعل والنتيجة مقبول ما دام لا يتجاوز $\tau_e$.
\end{proposition}
\end{keyblock}

\section{لماذا ينجح هذا: الضجيج بحثًا}
\label{sec:dofagradient}

لماذا تقود القاعدة~\eqref{eq:threefactor} إلى الأفضل؟ الجواب في الضجيج من الفصل~\ref{sec:code}. ليكن خرج العصبون $y=f(u)+\xi$، حيث $u=\sum_jw_jx_j$، و$\xi$ ارتعاش عشوائي متوسطه صفر وتباينه $\sigma^2$. المكافأة $R$ تعتمد على $y$. لنأخذ أثرًا لا يساوي $x_jy$ ببساطة، بل جزأه العشوائي $x_j\xi$، ولنأخذ عاملًا ثالثًا هو الخطأ $\delta=R-\bar R$، حيث $\bar R$ المكافأة المتوقعة. عند ضجيج صغير $R\approx R\bigl(f(u)\bigr)+R'\,\xi$، لذلك
\begin{empheq}[box=\keybox]{equation}
\bigl\langle \delta\,\xi\,x_j \bigr\rangle \;=\; \sigma^2\,R'\,x_j \;=\; \frac{\sigma^2}{f'(u)}\,\frac{\partial\langle R\rangle}{\partial w_j} .
\label{eq:perturbation}
\end{empheq}
الخطوة المتوسطة تسير على تدرج المكافأة المتوقعة~\cite{Williams1992}. لم يحسب أحد أي مشتقة: انحرفت الشبكة انحرافًا عشوائيًا، وأخبر الدوبامين هل صار الوضع أفضل، فتحركت المشابك المشاركة في الاتجاه المفيد (الشكل~\ref{fig:perturbation}). الضجيج الذي أعاق نقل الرسائل في الفصل~\ref{sec:code} يصير هنا بحثًا.

\begin{figure}[t]
\centering
\begin{tikzpicture}[>=Stealth,x=2cm,y=3cm]
\draw[->,gray!70!black] (0,0) -- (4.2,0) node[below left,font=\small,black] {الخرج $y$};
\draw[->,gray!70!black] (0,0) -- (0,1.2) node[left,font=\small,black] {المكافأة $R$};
% reward hill
\draw[very thick,blue!60!black] plot[domain=0:4,samples=60] (\x,{max(0,1-((\x-3)/2.2)^2)});
\node[font=\scriptsize,blue!60!black,anchor=south] at (3,1.02) {أفضل خرج};
% noise around f(u)
\draw[thick,gray!60!black] plot[domain=0.6:2.4,samples=50] (\x,{0.28*exp(-((\x-1.5)/0.3)^2/2)});
\draw[densely dashed,gray!60!black] (1.5,0) -- (1.5,0.535);
\node[below,font=\small] at (1.5,0) {$f(u)$};
\node[font=\scriptsize,gray!50!black] at (1.5,-0.2) {الارتعاش $\xi$};
% expected reward
\draw[densely dotted,gray!60!black] (0.5,0.535) node[left,font=\scriptsize,black] {$\bar R$} -- (2.1,0.535);
% two random deviations
\fill[green!45!black] (2.1,0.833) circle (0.035);
\draw[very thick,green!45!black] (2.1,0.535) -- (2.1,0.833);
\node[font=\scriptsize,green!45!black,anchor=west,align=left] at (2.18,0.6) {$\xi>0$: أفضل،\\ $\delta>0$};
\fill[red!70!black] (0.9,0.089) circle (0.035);
\draw[very thick,red!70!black] (0.9,0.535) -- (0.9,0.089);
\node[font=\scriptsize,red!70!black,anchor=east,align=right] at (0.83,0.27) {$\xi<0$: أسوأ،\\ $\delta<0$};
% resulting step
\draw[->,very thick,orange!85!black] (1.55,0.4) -- (1.95,0.4) node[right,font=\scriptsize,black] {خطوة};
\end{tikzpicture}
\caption{الضجيج بحثًا~\eqref{eq:perturbation}. خرج العصبون يرتعش حول $f(u)$. الانحراف العشوائي إلى اليمين (الأخضر) أعطى أكثر من المتوقع، $\delta>0$؛ والانحراف إلى اليسار (الأحمر) أعطى أقل، $\delta<0$. في الحالتين الجداء $\delta\,\xi$ موجب، فيتحرك الوزن إلى اليمين، حيث تزداد المكافأة. في المتوسط تتناسب الخطوة مع الميل $R'$: على قمة التلّ يكون الانحراف في الاتجاهين سيئًا بالقدر نفسه، فتلغي الخطوات بعضها بعضًا.}
\label{fig:perturbation}
\end{figure}

\begin{keyblock}
\begin{proposition}[النزول على التدرج بلا أسلاك عكسية]
\label{prop:gradient}
إذا كانت في الشبكة انحرافات عشوائية في النشاط، وكانت إشارة البثّ تساوي خطأ التنبؤ بالمكافأة ومتوسطها صفر في كل موقف، فإن القاعدة~\eqref{eq:threefactor} تغيّر الأوزان في المتوسط على تدرج المكافأة المتوقعة. ولا حاجة لذلك إلى أسلاك عكسية ولا إلى طور تعلّم مستقل.
\end{proposition}
\end{keyblock}

الآن يتضح لماذا يُخبر الدوبامين لا عن المكافأة، بل عن \emph{خطأ} التنبؤ بها. لو كان العامل الثالث المكافأةَ نفسها $R\ge0$ لما تغيّر المتوسط في~\eqref{eq:perturbation}، لكن لظهرت مشكلتان. الأولى أن الجزء المفيد سيُضاف إليه الحد $\bar R\,\xi x_j$: متوسطه صفر، لكنه ضجيج قوي. والثانية، وهي أسوأ، أن المشبك الحقيقي لا يستطيع أن يفصل في الأثر $x_jy$ الجزء العشوائي عن غير العشوائي. عند مداخل معينة يكون $x_jf(u)$ العدد نفسه من محاولة إلى أخرى؛ فإذا كان متوسط $\delta$ في هذا الموقف صفرًا فإن هذا الجزء من الأثر لا يغيّر شيئًا، أما عند $\delta\ge0$ فإنه يقوّي مرة بعد مرة كل ما فعلته الشبكة، صوابًا كان أم خطأً. لذلك يجب أن يكون $\bar R$ التوقعَ \emph{في الموقف المعين}، لا المتوسط على الحياة كلها. وحسابه مهمة الناقد الذي نتحدث عنه أدناه.

تحققنا من ذلك على نموذج. مجموعتان من عصبونات \nt{GLUT} تختاران أحد فعلين عبر التثبيط المتبادل، كما في الشكل~\ref{fig:wta}؛ الأوزان من إشارة ”ابدأ“ إلى المجموعتين متساوية في البداية، والضجيج يقرر من يفوز. الفعل $A$ يجلب العصير باحتمال $0.8$، والفعل $B$ باحتمال $0.2$، ويصل العصير بعد نصف ثانية من الاختيار. عند $\tau_e=1\,\text{s}$ و$\delta=R-\bar R$ ارتفعت نسبة اختيار $A$ على خمسمئة محاولة من $0.65\text{--}0.8$ في المئة الأولى إلى $0.96\text{--}1.0$ في الأخيرة، وبقيت الأوزان ضمن $0.85\text{--}1.35$. وعند $\tau_e=50\ms$، وهو أقل من تأخر المكافأة، لم تتعلم الشبكة شيئًا: بقيت نسبة اختيار $A$ نحو النصف. وعند $\delta=R$ دون طرح التوقع صارت الشبكة تختار $A$ أيضًا، لكن وزن الفائز ازداد خلال الخمسمئة محاولة نفسها ألف مرة واستمر في الازدياد؛ ومع $\delta=R$ نفسها ومعدل تعلّم أكبر عشر مرات ثبّتت الشبكة في واحدة من ثلاث محاكاة اختيارها العشوائي الأول إلى الأبد، وكان الأسوأ.

\section{ثمن السلك الواحد}

الإشارة $\delta$ واحدة للجميع، أما الضجيج الذي تقيّمه فهو مجموع ارتعاشات كل العصبونات $N$ المؤثرة في النتيجة. يرى كل مشبك في $\delta$ جزأه المفيد وضجيج الجيران الآخرين $N-1$. ولكي يستخلص نصيبه عليه أن يأخذ المتوسط على محاولات كثيرة، فيزداد زمن التعلّم متناسبًا مع $N$~\cite{Werfel2005}. أما الانتشار العكسي فلا يتطلب هذا الثمن.

\begin{keyblock}
\begin{proposition}[ثمن البثّ]
\label{prop:broadcastcost}
زمن التعلّم بإشارة مشتركة واحدة يزداد متناسبًا مع عدد العصبونات التي يؤثر ضجيجها في النتيجة. هذا التعلّم يصلح لدارات صغيرة تختار من بين خيارات قليلة، لا لضبط مليارات الأوزان.
\end{proposition}
\end{keyblock}

يبدو أن الدماغ يقسّم العمل هكذا بالضبط: تذهب مخارج عصبونات \nt{DOPA} قبل كل شيء إلى المناطق التي تختار الأفعال (الفصل~\ref{sec:neurontypes})، أما شبكة \nt{GLUT} الهائلة في القشرة، حيث الغالبية الساحقة من الأوزان، فتتعلم أساسًا بقواعد محلية، بلا معلّم خارجي. كانت هذه القواعد تُردّ سابقًا إلى قواعد هيب~\eqref{eq:oja}. أما الآن فتتزايد المعطيات على أن للقشرة هدفًا خاصًا بها، هو أن \emph{تتنبأ} بدخلها التالي، وعندئذ يُحسب الخطأ في مكانه، فرقًا بين المتنبَّأ به والواصل~\cite{Keller2018}: المعلّم مدمج في الشبكة نفسها. ونوافذ اللدونة التي طولها ثوانٍ، حيث تكون الإشارة الثالثة الإثارةَ القوية للمتلقي نفسه~\cite{Bittner2017}، تبيّن أن إشارة خطأ محلية تصل إلى المشابك فعلًا. أما إلى أي حد هذه قاعدة عامة في القشرة ففرضية.

\section{من أين يأتي الخطأ: الناقد في الزمن المتصل}

يبقى السؤال: كيف تحسب عصبونات \nt{DOPA} نفسها $\delta$؟ في برامج التعلّم المعزَّز يُحسب خطأ التنبؤ بخطوات $t,t+1,\dots$ وفي الكائن الحي لا توجد خطوات، لذلك نكتبه في الزمن المتصل~\cite{Doya2000}. ليكن $r(t)$ دفق المكافأة، و$V(t)$ توقع كل المكافأة المقبلة، حيث تُقدَّر المكافأة الأبعد بقيمة أقل:
\begin{equation}
V(t) \;=\; \int_t^{\infty} e^{-(s-t)/\tau_\gamma}\,r(s)\,\dd s .
\end{equation}
بالاشتقاق نحصل على $\dd V/\dd t=V/\tau_\gamma-r$. والفرق الذي يخالف به الواقعُ هذه المساواة هو خطأ التنبؤ:
\begin{empheq}[box=\keybox]{equation}
\delta(t) \;=\; r(t) \;+\; \frac{\dd V}{\dd t} \;-\; \frac{V}{\tau_\gamma} .
\label{eq:ctd}
\end{empheq}
الثابت $\tau_\gamma$ هو أفق التخطيط، أي النظير المتصل للمعامل $\gamma$؛ ووفق فرضيةٍ يضبطه \nt{SERO} (القسم~\ref{sec:horizon})، وعندئذ تكون~\eqref{eq:patience} قاعدةً لمقدار ما يستحق الانتظار. لنتحقق من ثلاث حالات (الشكل~\ref{fig:ctdcases}). أضاء مصباح يعد بالعصير: قفز $V$، و$\dd V/\dd t$ كبير، فتظهر ذروة. وصل العصير الموعود: $r>0$، لكن التوقع نفد في اللحظة نفسها، $\dd V/\dd t\approx-r$، فلا ذروة. لم يصل العصير: اختفى التوقع بلا مكافأة، $\dd V/\dd t<0$، فيحدث هبوط.

\begin{figure}[t]
\centering
\begin{tikzpicture}[>=Stealth,x=1.15cm,y=1cm]
\foreach \col/\ttl/\lampon/\juice in {0/{عصير غير متوقع}/0/1, 1/{عصير موعود}/1/1, 2/{لم يصل العصير}/1/0}{
 \begin{scope}[xshift=\col*4.6cm]
  \node[font=\small] at (1.5,3.55) {\ttl};
  \foreach \yy in {2.4,1.2,0}{\draw[gray!60!black] (0,\yy) -- (3.1,\yy);}
  % reward r
  \ifnum\juice=1
    \fill[green!45!black] (1.95,2.4) rectangle (2.05,2.85);
    \pic[scale=0.55] at (2.0,3.1) {juice};
  \else
    \pic[scale=0.55] at (2.0,3.1) {nojuice};
  \fi
  % expectation V and error delta
  \ifnum\lampon=1
    \pic[scale=0.55] at (1.0,3.1) {lamp};
    \draw[very thick,blue!60!black] (0,1.2) -- (1,1.2) -- (1,1.45)
      plot[domain=1:2,samples=20] (\x,{1.2+0.25*exp((\x-1)/1.5)}) -- (2,{1.2+0.25*exp(1/1.5)}) -- (2,1.2) -- (3.1,1.2);
    \draw[very thick,red!70!black] (0,0) -- (0.95,0) -- (0.95,0.5) -- (1.05,0.5) -- (1.05,0) -- (3.1,0);
    \ifnum\juice=0
      \draw[very thick,red!70!black] (1.95,0) -- (1.95,-0.45) -- (2.05,-0.45) -- (2.05,0);
      \node[font=\scriptsize,red!70!black,anchor=west] at (2.1,-0.35) {هبوط};
    \else
      \node[font=\scriptsize,gray!50!black,anchor=north,align=center] at (2.0,-0.08) {$r$ وانخفاض $V$\\ يلغي أحدهما الآخر};
    \fi
  \else
    \draw[very thick,blue!60!black] (0,1.2) -- (3.1,1.2);
    \draw[very thick,red!70!black] (0,0) -- (1.95,0) -- (1.95,0.5) -- (2.05,0.5) -- (2.05,0) -- (3.1,0);
  \fi
  \ifnum\col=0
    \node[font=\small,anchor=east] at (-0.1,2.55) {$r$};
    \node[font=\small,anchor=east] at (-0.1,1.35) {$V$};
    \node[font=\small,anchor=east] at (-0.1,0.1) {$\delta$};
  \fi
 \end{scope}}
\end{tikzpicture}
\caption{ثلاث حالات للخطأ~\eqref{eq:ctd}؛ من الأعلى إلى الأسفل: المكافأة $r$، والتوقع $V$، والخطأ $\delta$. على اليسار لا يوجد توقع، فالعصير مفاجأة كله. في الوسط يرفع المصباح $V$ قفزةً واحدة: هذه قفزة في $\dd V/\dd t$، فتقع ذروة $\delta$ عند المصباح. بعد ذلك ينمو $V$ تمامًا كما يقتضي الأفق، و$\delta=0$؛ ولحظة العصير تلغي المكافأة $r$ وانخفاض $V$ أحدهما الآخر. على اليمين ينخفض $V$ بلا مكافأة، فيحدث هبوط. هذه هي الذروة والانتقال والهبوط أنفسها التي في الشكل~\ref{fig:rpe}، لكننا نرى الآن من أين تأتي.}
\label{fig:ctdcases}
\end{figure}

ماذا يلزم لـ~\eqref{eq:ctd} من الدارات التي لدينا؟ ثلاثة أشياء.

\emph{التقدير $V$ نفسه.} تخرجه مجموعة من عصبونات \nt{GLUT}، هي الناقد، تتعلم أوزانها بالقاعدة نفسها~\eqref{eq:threefactor} وبـ$\delta$ نفسها: إذا كان $\delta>0$ فالتوقع كان أقل من اللازم، فتقوى الوصلات التي كانت نشطة قبل ذلك.

\emph{المشتقة $\dd V/\dd t$.} تحسبها أي دارة تستجيب للتغيرات لا للمستوى: إثارة مباشرة مع تثبيط متأخر عبر وسيط \nt{GABA}، كما في كاشف الاتجاه في الفصل~\ref{sec:gaba}، أو مشبك مستنزَف~\eqref{eq:depression} من الفصل~\ref{sec:code}.

\emph{الإشارة في سلك واحد.} الخطأ قد يكون موجبًا أو سالبًا، أما تردد النبضات فموجب فقط. والحل هو نفسه عند عصبون \nt{SERO} في الفصل~\ref{sec:serocomputer}: عصبون \nt{DOPA} ناظمة إيقاع. نبضاته المنتظمة القليلة في الثانية تعني $\delta=0$، والرشقة تعني $\delta>0$، والتوقف يعني $\delta<0$. ويبقى للخطأ السالب مجال أضيق: لا يمكن خفض التردد تحت الصفر. وعند عصبونات \nt{DOPA} الحقيقية تكون الهبوطات فعلًا أضحل من الذرى~\cite{Bayer2005}.

أن الدوبامين يسلك سلوك $\delta$ أمر مقيس. أما كيف يحسب الدماغ $\dd V/\dd t$ بالضبط ففرضية.

\section{الفاعل والناقد}

لنجمع كل شيء (الشكل~\ref{fig:actorcritic}). عصبونات الموقف تثير مجموعات الأفعال التي تختار الفائز عبر \nt{GABA} (القضية~\ref{prop:wta})، وهذا هو \emph{الفاعل}، وتثير كذلك الناقد الذي يخرج $V$~\cite{SuttonBarto2018}. يتلقى عصبون \nt{DOPA} المكافأة $r$ و$V$، ويحسب~\eqref{eq:ctd}، ويبثّ $\delta$ إلى كل مشابك الفاعل والناقد.

\begin{figure}[t]
\centering
\begin{tikzpicture}[>=Stealth,
  nrn/.style={circle,draw,very thick,fill=blue!6,minimum size=0.9cm,inner sep=0pt,font=\small},
  gab/.style={circle,draw,very thick,fill=red!10,minimum size=0.6cm,inner sep=0pt},
  dofa/.style={circle,draw,very thick,double,double distance=1.2pt,fill=orange!15,minimum size=1.35cm,inner sep=0pt,font=\scriptsize},
  inh/.style={-{Bar[width=7pt]},thick,red!70!black},
  bc/.style={->,thick,densely dashed,orange!85!black}]
\node[nrn] (S1) at (0,2.4) {$x_1$};
\node[nrn] (S2) at (0,1.2) {$x_2$};
\node[nrn] (S3) at (0,0) {$x_3$};
\node[left=0.15cm,align=right,font=\small] at (-0.5,1.2) {الموقف};
\node[nrn] (A) at (4,2.6) {$A$};
\node[nrn] (B) at (4,1.0) {$B$};
\node[gab] (G) at (5.2,1.8) {};
\draw[->,thick] (A) -- (G); \draw[->,thick] (B) -- (G);
\draw[inh] (G) to[bend left=35] (A);
\draw[inh] (G) to[bend right=35] (B);
\node[above,font=\small] at (4.6,3.05) {الفاعل: اختيار الفعل};
\node[nrn] (V) at (4,-1.1) {$V$};
\node[below,font=\small] at (4,-1.6) {الناقد};
\foreach \s in {S1,S2,S3}{\foreach \t in {A,B,V}{\draw[->,gray!60!black] (\s) -- (\t);}}
\node[dofa] (D) at (8.0,-1.1) {\nt{DOPA}};
\draw[->,thick] (V) -- node[above,font=\scriptsize] {$V,\ \dd V/\dd t$} (D);
\draw[->,thick,green!45!black] (10.0,-1.1) node[right,black,font=\small] {المكافأة $r$} -- (D);
\pic[scale=1.1] at (10.6,-0.5) {juice};
\draw[bc] (D.north) -- (8.0,4.0) -- node[above,font=\small,black] {$\delta$ إلى كل مشابك الفاعل والناقد} (2.2,4.0) -- (2.2,3.0);
\draw[bc] (D.south) -- (8.0,-2.4) -- (2.2,-2.4) -- (2.2,-0.6);
\end{tikzpicture}
\caption{شبكة تتعلم اختيار الأفعال. الأسهم الرمادية مشابك \nt{GLUT} بأوزان متغيرة؛ الدائرة الحمراء عصبون \nt{GABA}؛ الدائرة المزدوجة ناظمة إيقاع \nt{DOPA} تحسب $\delta$ وفق~\eqref{eq:ctd}؛ الأسهم المتقطعة بثّ $\delta$.}
\label{fig:actorcritic}
\end{figure}

إشارة أثر الناقل يختارها المتلقي (القضية~\ref{prop:receptor})، وفي منطقة اختيار الأفعال يحمل جزء من العصبونات مستقبلات دوبامين من عائلة، وجزء آخر من عائلة أخرى. في تجارب على الشرائح يقوّي الدوبامين مع النشاط المشترك المشابكَ عند الأولى، ويضعفها عند الثانية~\cite{Shen2008}؛ وفي النموذج يعني هذا أن إشارة $\delta$ في~\eqref{eq:threefactor} معكوسة عند الثانية. الأولى تتعلم ما \emph{يستحق الفعل}، والثانية ما \emph{لا يستحق الفعل}.

\section{الخلاصة}

\begin{table}[t]
\centering
\footnotesize
\setlength{\tabcolsep}{4pt}
\renewcommand{\arraystretch}{1.15}
\begin{tabular}{>{\raggedright}p{2.6cm}>{\raggedright}p{2.7cm}>{\raggedright}p{3.1cm}>{\raggedright\arraybackslash}p{5.0cm}}
\toprule
القاعدة & ما يعرفه المشبك & ما تتعلمه الشبكة & ما هو معروف \\
\midrule
هيب بالترددات~\eqref{eq:oja} & $x$، $y$، وزنه & ضغط المداخل: الاتجاهات الرئيسية & التقوية عند النشاط المشترك مقيسة؛ وآلية تقييد الأوزان موضوع بحث \\
هيب بالتوقيت & ترتيب نبضات $x$ و$y$ ضمن $\sim20\ms$ & الروابط السببية، والتسلسلات & مقيسة على مشابك \nt{GLUT} \\
العوامل الثلاثة~\eqref{eq:threefactor} & $x$، $y$، الأثر $e$، الإشارة المشتركة $\delta$ & الأفعال التي تجلب المكافأة & تأثير الدوبامين خلال ثوانٍ بعد النشاط مقيس؛ وكونها الآلية الرئيسية لتعلّم الاختيار فرضية قوية الأساس \\
الناقد~\eqref{eq:ctd} & كذلك & توقع المكافأة $V$ & تشابه الدوبامين مع $\delta$ مقيس؛ والدارة التي تحسب $\dd V/\dd t$ فرضية \\
الانتشار العكسي للخطأ & تصحيحًا خاصًا لكل وزن & كل شيء، لكنه يحتاج إلى أسلاك عكسية وساعة & لم يوجد في الدماغ بهذا الشكل؛ ويُبحث عن تقريبات له \\
\bottomrule
\end{tabular}
\caption{قواعد التعلّم في شبكة من عصبونات \nt{GLUT} و\nt{GABA} و\nt{DOPA}.}
\label{tab:learning}
\end{table}

قواعد التعلّم مجموعة في الجدول~\ref{tab:learning}. \nt{GLUT} يحمل البيانات، و\nt{GABA} يدير التدفق، و\nt{DOPA} يقرر أي التغييرات في الشبكة تبقى.

\section{تمارين}

\begin{exercise}[\lvlA]
\label{ex:06:1}
\emph{كم يعيش الأثر.} يدوم النشاط المشترك $xy=1$ مدة $T_a=0.2\,\text{s}$ بدءًا من $e=0$، ويصل الدوبامين بعد $d=0.5\,\text{s}$ من نهايته. (أ)~كم مرة يكون الأثر~\eqref{eq:threefactor} في تلك اللحظة أكبر عند $\tau_e=1\,\text{s}$ منه عند $\tau_e=50\ms$؟ (ب)~عند أي $\tau_e$ يكون أكبر ما يمكن؟
\end{exercise}

\begin{exercise}[\lvlA]
\label{ex:06:2}
\emph{انجراف قاعدة هيب.} المرسِل والمتلقي يخرجان دفقين بواسونيين مستقلين بمعدل $10$ نبضات في الثانية، وكل زوج من النبضات يغيّر الوزن وفق نافذة الشكل~\ref{fig:stdp} مع $\tau_\pm=20\ms$ و$A_-=0.6A_+$. بكم $A_+$ يزداد الوزن خلال ساعة من نشاط عشوائي تمامًا؟ عند أي $\tau_-$ يختفي الانجراف؟
\end{exercise}

\begin{exercise}[\lvlB]
\label{ex:06:3}
\emph{ارتباطات لا تغايرات.} قاعدة القضية~\ref{prop:oja} تجد المتجه الذاتي الرئيسي لـ$C=\langle\mathbf x\mathbf x^{\mathsf T}\rangle$. الترددات موجبة: $\mathbf x=\mathbf m+\mathbf n$، $\mathbf m=(\mu,0)$، وتغاير الضجيج $\operatorname{diag}(s_1^2,s_2^2)$. بأي شرط يتعلم العصبون $\mathbf e_1$؟ وماذا يتعلم عند $\mu=1$، $s_1=0.1$، $s_2=0.5$؟
\end{exercise}

\begin{exercise}[\lvlA]
\label{ex:06:4}
\emph{الأفق في الذروة.} يضيء مصباح في لحظة لا يمكن التنبؤ بها، ويصل عصير حجمه $R$ بعد $L=1\,\text{s}$. بيّن أنه للتوقع المتعلَّم وفق~\eqref{eq:ctd} يكون $\delta\equiv0$ بين المصباح والعصير، وجد مساحة الذروة عند المصباح عند $\tau_\gamma=2\,\text{s}$ و$\tau_\gamma=4\,\text{s}$.
\end{exercise}

\begin{exercise}[\lvlB]
\label{ex:06:5}
\emph{من أين يأتي ثمن البثّ.} $N$ عصبونًا ترتعش باستقلال، $\xi_i\sim\mathcal N(0,\sigma^2)$، والخطأ $\delta=a\sum_i\xi_i$، ويقدّر المشبك $j$ التدرج بـ$\delta\,\xi_jx_j$. جد نسبة الإشارة إلى الضجيج في محاولة واحدة. عصبون واحد يتعلم في $100$ محاولة مدة كل منها $5\,\text{s}$: كم يستغرق التعلّم عند $N=10^4$ و$N=10^{10}$؟
\end{exercise}

\begin{exercise}[\lvlB]
\label{ex:06:6}
\emph{تصميم: دارة تحسب $\delta$.} يتلقى عصبون \nt{DOPA} المقدار $r$، و$V$ بوزن $k_1$، ونسخة من $V$ منعَّمة بـ$\tau_d$ بوزن $-k_2$. (أ)~اختر $k_1$ و$k_2$ بحيث يكون الخرج للمقادير $V$ البطيئة هو~\eqref{eq:ctd}. (ب)~عند $V=V_0e^{t/\tau_\gamma}$ تكون $\delta$ الحقيقية صفرًا. عند أي $\tau_d$ لا يتجاوز الخطأ الزائف للدارة $5\,\%$ من $V/\tau_\gamma$ إذا كان $\tau_\gamma=2\,\text{s}$؟
\end{exercise}

\begin{exercise}[\lvlB]
\label{ex:06:7}
\emph{نمذجة: أقصى تأخر للمكافأة.} أضف إلى \texttt{run()} في نسخة من \texttt{models/\allowbreak dofa\_learning.py} تأخرًا للمكافأة $d$. جد نسبة اختيار $A$ في المئة الأخيرة من $500$ محاولة عند $d=0.5\,\text{s}$ و$\tau_e=0.05$ و$0.2$ و$1$ و$4\,\text{s}$، وعند $\tau_e=1\,\text{s}$ و$d=3\,\text{s}$. عند أي نسبة من الأثر الباقي حتى المكافأة (التمرين~\ref{ex:06:1}) يختفي التعلّم؟
\end{exercise}

\begin{exercise}[\lvlC]
\label{ex:06:8}
\emph{هل يمكن تفادي ثمن البثّ؟} $N$ عصبونًا: $y_i=w_i+\xi_i$، $\xi_i\sim\mathcal N(0,\sigma^2)$، والمكافأة $R=-\sum_i(y_i-w_i^*)^2$، والقاعدة $\Delta w_i=\eta\,(R-\langle R\rangle)\,\xi_i$. في كم محاولة، عند أفضل $\eta$، يتناقص الخطأ $\sum_i(w_i-w_i^*)^2$ بمعامل $e$؟ وماذا لو ارتعشت $m$ عصبونات عشوائية فقط من $N$؟ هل يساعد البحث المتناثر؟
\end{exercise}
